\section{Memulai: Menjalankan JavaScript di Browser}

JavaScript dapat dijalankan langsung di browser melalui tiga cara: inline (di dalam atribut event HTML), internal (dalam tag \code{<script>} di file HTML), dan eksternal (file \code{.js} terpisah yang dihubungkan dengan \code{<script src="...">}). Cara eksternal adalah yang paling direkomendasikan karena memisahkan logika dari struktur, memudahkan caching, dan meningkatkan maintainability \cite{mdnJsGuide}.

\begin{table}[htbp]
\centering
\begin{tabular}{|P{2.5cm}|P{4.5cm}|P{5cm}|}
\hline
\textbf{Metode} & \textbf{Cara} & \textbf{Rekomendasi} \\
\hline
Inline & \code{onclick="alert('Hi')"} & Hindari; sulit dirawat \\
\hline
Internal & \code{<script>} di HTML & Untuk skrip kecil/demo \\
\hline
Eksternal & \code{<script src="app.js">} & Direkomendasikan; pisahkan file \\
\hline
\end{tabular}
\caption{Tiga cara menyisipkan JavaScript di HTML}
\end{table}

Browser menyediakan \textbf{DevTools Console} (dibuka dengan F12 atau Ctrl+Shift+I) yang memungkinkan menulis dan menjalankan JavaScript secara interaktif. Console sangat berguna untuk debugging: \code{console.log()} menulis pesan ke console, \code{console.error()} untuk error, \code{console.table()} menampilkan data tabular. Setiap browser modern (Chrome, Firefox, Edge, Safari) memiliki DevTools dengan fitur serupa \cite{jsInfo}.

\begin{catatan}
\textbf{Atribut \code{defer} dan \code{async}.} Tag \code{<script>} secara default memblokir parsing HTML saat dieksekusi. Atribut \code{defer} menunda eksekusi sampai HTML selesai diparsing (urutan dipertahankan). Atribut \code{async} menjalankan skrip segera setelah diunduh (urutan tidak dijamin). Untuk skrip yang memanipulasi DOM, selalu gunakan \code{defer} atau letakkan \code{<script>} sebelum \code{</body>} \cite{mdnJsGuide}.
\end{catatan}

